\section{Variables, loops, lists, and arrays}

\begin{frame}[fragile]{Do you have access to Python?}
  For this course we use \href{https://docs.astral.sh/uv/}{uv} to install Python and packages.

  \begin{enumerate}
    \item Install \texttt{uv}
    \item At the project root (folder with \texttt{pyproject.toml}):
  \end{enumerate}
\begin{lstlisting}[style=terminal]
uv sync
cd python_intro
uv run python distance.py
\end{lstlisting}
  \texttt{uv run} uses the project environment; no manual \texttt{venv} activation needed.
\end{frame}

\begin{frame}{Mathematical example}
  Most examples will involve this formula:
  \begin{equation}
    s = v_0 t + \frac{1}{2} a t^2
  \end{equation}
  We may view $s$ as a function of $t$: $s(t)$, and also include the
  parameters in the notation: $s(t;v_0,a)$.
\end{frame}

\begin{frame}[fragile]{A program for evaluating a formula}
  \textbf{Task:} Compute $s$ for $t=0.5$, $v_0=2$, and $a=0.2$.

  \textbf{Python code} (\texttt{python\_intro/distance.py}):
\lstinputlisting[style=pythoncompact,firstline=6]{code/distance.py}
\end{frame}

\begin{frame}[fragile]{Execution}
\begin{lstlisting}[style=terminal]
Terminal> uv run python distance.py
1.025
s=1.025
s        =       1.025
\end{lstlisting}
\end{frame}

\begin{frame}[fragile]{Assignment statements assign a name to an object}
\begin{lstlisting}[style=pythoncompact]
t = 0.5                   # real number makes float object
v0 = 2                    # integer makes int object
a = 0.2                   # float object
s = v0*t + 0.5*a*t**2     # float object
\end{lstlisting}
  Rule:
  \begin{itemize}
    \item evaluate right-hand side; it results in an \emph{object}
    \item left-hand side is a name for that object
  \end{itemize}
\end{frame}

\begin{frame}[fragile]{Formatted output with text and numbers}
  \begin{itemize}
    \item Task: write out text with a number (3 decimals): \texttt{s=1.025}
    \item Method: f-strings (modern Python)
  \end{itemize}
\begin{lstlisting}[style=pythoncompact]
print(f's={s:g}')       # g: compact notation
print(f's={s:.2f}')     # f: decimal notation, 2 decimals
\end{lstlisting}
  Older alternatives (still valid):
\begin{lstlisting}[style=pythoncompact]
print('s=%g' % s)
print('s={s:.2f}'.format(s=s))
\end{lstlisting}
\end{frame}

\begin{frame}[fragile]{Programming with a while loop}
  \begin{itemize}
    \item Task: write out a table of $t$ and $s(t)$ values (two columns),
          for $t\in [0,2]$ in steps of $0.1$
    \item Method: while loop
  \end{itemize}
\lstinputlisting[style=pythonTiny,firstline=7]{code/while.py}
\end{frame}

\begin{frame}[fragile]{Output of the previous program}
\begin{lstlisting}[style=terminal]
Terminal> uv run python while.py
0 0.0
0.1 0.201
0.2 0.404
0.3 0.609
0.4 0.816
0.5 1.025
...
1.8 3.924
1.9 4.161
\end{lstlisting}
\end{frame}

\begin{frame}[fragile]{Structure of a while loop}
\begin{lstlisting}[style=pythoncompact]
while condition:
    <indented statement>
    <indented statement>
    <indented statement>
\end{lstlisting}
  Note:
  \begin{itemize}
    \item the colon in the first line
    \item all statements in the loop \emph{must be indented}\\
          (no braces as in C, C++, Java, \ldots)
    \item \texttt{condition} is a boolean expression (e.g., \texttt{t <= 2})
  \end{itemize}
\end{frame}

\begin{frame}[fragile]{Closer look at the output}
\begin{lstlisting}[style=terminal]
Terminal> uv run python while.py
0 0.0
0.1 0.201
...
1.8 3.924
1.9 4.161
\end{lstlisting}
  The last line contains $1.9$, but the while loop should run also when
  $t=2$ since the test is \texttt{t <= 2}. Why is this test \texttt{False}?
\end{frame}

\begin{frame}[fragile]{Examine the program (Python Tutor idea)}
  Step through a smaller loop and examine variables
  (\href{https://pythontutor.com/visualize.html\#mode=edit}{Python Tutor}):
\begin{lstlisting}[style=pythoncompact]
a = 0
da = 0.4
while a <= 1.2:
    print(a)
    a = a + da
\end{lstlisting}
\end{frame}

\begin{frame}[fragile]{Round-off errors!}
  Why is \texttt{a <= 1.2} false when \texttt{a} ``is'' $1.2$?
\begin{lstlisting}[style=pythoncompact]
a = 0
da = 0.4
while a <= 1.2:
    print(a)
    a = a + da
    print(f'da={da:.16E}\na={a:.16E}')
    print(f'a <= 1.2: {a <= 1.2}')
\end{lstlisting}
  \begin{alertblock}{Key point}
    Small round-off error in \texttt{da} makes
    \texttt{a = 1.2000000000000002}
  \end{alertblock}
\end{frame}

\begin{frame}[fragile]{Never \texttt{a == b} for real numbers}
  Always use a tolerance!
\begin{lstlisting}[style=pythoncompact]
a = 1.2
b = 0.4 + 0.4 + 0.4
boolean_condition1 = a == b              # may be False

# This is the way to do it
tol = 1E-14
boolean_condition2 = abs(a - b) < tol    # True
\end{lstlisting}
\end{frame}

\begin{frame}[fragile]{A list collects several objects in a sequence}
  A list of numbers:
\begin{lstlisting}[style=pythoncompact]
L = [-1, 1, 8.0]
\end{lstlisting}
  A list can contain any type of objects:
\begin{lstlisting}[style=pythoncompact]
L = ['mydata.txt', 3.14, 10]   # string, float, int
\end{lstlisting}
\end{frame}

\begin{frame}[fragile]{Basic list operations}
\begin{lstlisting}[style=pythonTiny]
>>> L = ['mydata.txt', 3.14, 10]
>>> print(L[0])    # first element (index 0)
mydata.txt
>>> print(L[1])    # second element (index 1)
3.14
>>> del L[0]       # delete the first element
>>> print(L)
[3.14, 10]
>>> print(len(L))  # length of L
2
>>> L.append(-1)   # add -1 at the end
>>> print(L)
[3.14, 10, -1]
\end{lstlisting}
\end{frame}

\begin{frame}[fragile]{Store our table in two lists}
\lstinputlisting[style=pythonTiny,firstline=7]{code/while_list.py}
\end{frame}

\begin{frame}[fragile]{For loops}
  A for loop visits elements in a list, one by one:
\begin{lstlisting}[style=pythoncompact]
>>> L = [1, 4, 8, 9]
>>> for e in L:
...     print(e)
...
1
4
8
9
\end{lstlisting}
\end{frame}

\begin{frame}[fragile]{For-loop demo (Python Tutor style)}
\begin{lstlisting}[style=pythoncompact]
list1 = [0, 0.1, 0.2]
list2 = []
for element in list1:
    p = element + 2
    list2.append(p)
print(list2)
\end{lstlisting}
\end{frame}

\begin{frame}[fragile]{Integer counter over list/array indices}
\begin{lstlisting}[style=pythoncompact]
somelist = ['file1.dat', 22, -1.5]

for i in range(len(somelist)):
    # access list element through index
    print(somelist[i])
\end{lstlisting}
  Note:
  \begin{itemize}
    \item \texttt{range} returns a sequence of integers
    \item \texttt{range(a, b, s)} gives \texttt{a, a+s, \ldots} up to
          \emph{but not including} $b$
    \item \texttt{range(b)} implies $a=0$ and $s=1$
    \item \texttt{range(len(somelist))} gives $0,1,2$
  \end{itemize}
\end{frame}

\begin{frame}[fragile]{Replace our while loop by a for loop}
\lstinputlisting[style=pythonTiny,firstline=7]{code/for_list.py}
\end{frame}

\begin{frame}[fragile]{Traversal of multiple lists with \texttt{zip}}
\begin{lstlisting}[style=pythoncompact]
for e1, e2, e3, ... in zip(list1, list2, list3, ...):
\end{lstlisting}
  Alternative: loop over a common index
\begin{lstlisting}[style=pythoncompact]
for i in range(len(list1)):
    e1 = list1[i]
    e2 = list2[i]
    e3 = list3[i]
\end{lstlisting}
\end{frame}

\begin{frame}{Arrays are efficient lists of numbers}
  \begin{itemize}
    \item Lists collect objects in a single variable
    \item Lists are flexible (can grow, can contain ``anything'')
    \item Array: computationally efficient list of numbers
    \item Arrays have fixed length and one numeric type
    \item Arrays require the \texttt{numpy} module
  \end{itemize}
\end{frame}

\begin{frame}[fragile]{Examples on using arrays}
\begin{lstlisting}[style=pythonTiny]
>>> import numpy
>>> L = [1, 4, 10.0]    # List of numbers
>>> a = numpy.array(L)  # Convert to array
>>> print(a)
[ 1.  4. 10.]
>>> print(a[1])         # indexing
4.0
>>> print(a[0:2])       # slice (index 2 not included!)
[1. 4.]
>>> print(a.dtype)      # data type of an element
float64
>>> b = 2*a + 1         # arithmetics on arrays
>>> print(b)
[ 3.  9. 21.]
\end{lstlisting}
\end{frame}

\begin{frame}[fragile]{\texttt{numpy} functions create entire arrays}
  Apply $\ln$ to all elements in array \texttt{a}:
\begin{lstlisting}[style=pythoncompact]
>>> c = numpy.log(a)
>>> print(c)
[0.         1.38629436 2.30258509]
\end{lstlisting}
  Create $n+1$ uniformly spaced coordinates in $[a,b]$:
\begin{lstlisting}[style=pythoncompact]
t = numpy.linspace(a, b, n+1)
\end{lstlisting}
  Array of length $n$ filled with zeros:
\begin{lstlisting}[style=pythoncompact]
t = numpy.zeros(n)
s = numpy.zeros_like(t)  # zeros with t's size and dtype
\end{lstlisting}
\end{frame}

\begin{frame}[fragile]{Let's use arrays in our previous program}
\lstinputlisting[style=pythonTiny,firstline=7]{code/for_array.py}
  Note: no explicit loop for computing \texttt{s\_values}!
\end{frame}

\begin{frame}[fragile]{Standard math functions: the \texttt{math} module}
\begin{lstlisting}[style=pythoncompact]
>>> import math
>>> print(math.sin(math.pi))
1.2246467991473532e-16         # approximate value
\end{lstlisting}
  Get rid of the \texttt{math} prefix:
\begin{lstlisting}[style=pythoncompact]
from math import sin, pi
print(sin(pi))

from math import *   # import everything
print(sin(pi), log(e), tanh(0.5))
\end{lstlisting}
\end{frame}

\begin{frame}[fragile]{Use \texttt{numpy} for math on arrays}
  Matlab-style:
\begin{lstlisting}[style=pythoncompact]
from numpy import *
x = linspace(0, 1, 101)
y = exp(-x)*sin(pi*x)
\end{lstlisting}
  Preferred style:
\begin{lstlisting}[style=pythoncompact]
import numpy as np
x = np.linspace(0, 1, 101)
y = np.exp(-x)*np.sin(np.pi*x)
\end{lstlisting}
\end{frame}

\begin{frame}[fragile]{Our convention: \texttt{np} prefix, clean formulas}
\begin{lstlisting}[style=pythoncompact]
import numpy as np
x = np.linspace(0, 1, 101)

from numpy import sin, exp, pi
y = exp(-x)*sin(pi*x)
\end{lstlisting}
\end{frame}

\begin{frame}[fragile]{Array assignment gives a view (no copy!)}
  Consider array assignment \texttt{b = a}:
\begin{lstlisting}[style=pythoncompact]
a = np.linspace(1, 5, 5)
b = a
\end{lstlisting}
  Here, \texttt{b} is a \emph{view} (pointer) to the data of \texttt{a} ---
  no copying of data!
\begin{lstlisting}[style=pythonTiny]
>>> a
array([1., 2., 3., 4., 5.])
>>> b[0] = 5                          # changes a[0] to 5
>>> a
array([5., 2., 3., 4., 5.])
>>> a[1] = 9                          # changes b[1] to 9
>>> b
array([5., 9., 3., 4., 5.])
\end{lstlisting}
\end{frame}

\begin{frame}[fragile]{Copying array data: the \texttt{copy} method}
\begin{lstlisting}[style=pythonTiny]
>>> c = a.copy()       # copy all elements to new array c
>>> c[0] = 6           # a is not changed
>>> a
array([1., 2., 3., 4., 5.])
>>> c
array([6., 2., 3., 4., 5.])
\end{lstlisting}
\end{frame}

\begin{frame}{Construction of tridiagonal and sparse matrices}
  \begin{itemize}
    \item SciPy offers
          \href{https://docs.scipy.org/doc/scipy/reference/sparse.html}{scipy.sparse}
    \item Build sparse matrices from diagonals with \texttt{diags}
          (or the newer \texttt{diags\_array})
    \item Super-/sub-diagonals are shorter than the main diagonal
    \item Scalar broadcasting is supported (specify \texttt{shape})
  \end{itemize}
\end{frame}

\begin{frame}[fragile]{Tridiagonal matrix with \texttt{diags}}
\begin{lstlisting}[style=pythonTiny]
>>> import numpy as np
>>> import scipy.sparse
>>> N = 6
>>> diagonals = [
...     -np.linspace(1, N, N)[:-1],
...     -2*np.ones(N),
...     np.linspace(1, N, N)[1:],
... ]
>>> A = scipy.sparse.diags(diagonals, [-1, 0, 1], format='csc')
>>> A.toarray()
[[-2.  2.  0.  0.  0.  0.]
 [-1. -2.  3.  0.  0.  0.]
 [ 0. -2. -2.  4.  0.  0.]
 [ 0.  0. -3. -2.  5.  0.]
 [ 0.  0.  0. -4. -2.  6.]
 [ 0.  0.  0.  0. -5. -2.]]
\end{lstlisting}
\end{frame}

\begin{frame}[fragile]{Scalar broadcasting with \texttt{diags}}
\begin{lstlisting}[style=pythonTiny]
>>> B = scipy.sparse.diags(
...     [1, 2, 3], [-2, 0, 1], shape=(6, 6), format='csc')
>>> B.toarray()
[[2. 3. 0. 0. 0. 0.]
 [0. 2. 3. 0. 0. 0.]
 [1. 0. 2. 3. 0. 0.]
 [0. 1. 0. 2. 3. 0.]
 [0. 0. 1. 0. 2. 3.]
 [0. 0. 0. 1. 0. 2.]]
\end{lstlisting}
\end{frame}

\begin{frame}[fragile]{Legacy note: \texttt{spdiags}}
  Older code may use \texttt{spdiags}: all diagonals must have length $N$;
  unused entries sit at the start of super-diagonals and the end of
  sub-diagonals.
\begin{lstlisting}[style=pythonTiny]
>>> diagonals = np.zeros((3, N))
>>> diagonals[0, :] = np.linspace(-1, -N, N)
>>> diagonals[1, :] = -2
>>> diagonals[2, :] = np.linspace(1, N, N)
>>> A = scipy.sparse.spdiags(
...     diagonals, [-1, 0, 1], N, N, format='csc')
\end{lstlisting}
  Prefer \texttt{diags} / \texttt{diags\_array} in new code.
\end{frame}

\begin{frame}[fragile]{Solving a tridiagonal system}
  Solve $Ax=b$: choose $x$, compute $b=Ax$, solve, check.
\begin{lstlisting}[style=pythonTiny]
>>> x = np.linspace(-1, 1, N)
>>> b = A @ x                    # sparse matrix-vector product
>>> import scipy.sparse.linalg
>>> x = scipy.sparse.linalg.spsolve(A, b)
>>> print(x)
[-1.  -0.6 -0.2  0.2  0.6  1. ]
\end{lstlisting}
\end{frame}

\begin{frame}[fragile]{Check against dense matrix computations}
\begin{lstlisting}[style=pythonTiny]
>>> A_d = A.toarray()
>>> b = A_d @ x                  # dense product
>>> x = np.linalg.solve(A_d, b)
>>> print(x)
[-1.  -0.6 -0.2  0.2  0.6  1. ]
\end{lstlisting}
  See also \texttt{scipy\_intro/tridiagonal.py} in this repository.
\end{frame}

\begin{frame}[fragile]{Plotting}
  Plotting is done with \texttt{matplotlib}:
\lstinputlisting[style=pythonTiny,firstline=7]{code/plot1.py}
\end{frame}

\begin{frame}[fragile]{Plotting of multiple curves}
  Core idea (see \texttt{python\_intro/plot2.py}):
\begin{lstlisting}[style=pythonTiny]
v0, a0, a1, n = 0.2, 2, 3, 21
t = np.linspace(0, 2, n + 1)
s0 = v0 * t + 0.5 * a0 * t**2
s1 = v0 * t + 0.5 * a1 * t**2
plt.plot(t, s0, 'r-', t, s1, 'bo')
plt.xlabel('t [s]'); plt.ylabel('s [m]')
plt.legend([rf'$s(t; v_0={v0}, a={a0})$',
            rf'$s(t; v_0={v0}, a={a1})$'])
\end{lstlisting}
\end{frame}
